\section{附录D　内部资料对应与复现索引（用于核对定义、计数与实验接口）}

\begin{quote}
本附录仅给出与本文形式对象对应的内部资料定位，便于复现实验与核对计数；正文不依赖外部引文链条。
\end{quote}

\begin{enumerate}
  \item \textbf{Weyl 对与无公共本征向量：}给出 $U_{\mathrm{scan}}V=e^{2\pi i\alpha}VU_{\mathrm{scan}}$ 且 $\alpha\notin\QQ$ 时无公共本征向量的命题与证明。
  \item \textbf{黄金分支的 Hurwitz 极值与抗锁定解释：}黄金斜率的最坏逼近性质及其对相位锁定的延迟作用。
  \item \textbf{Sturmian 读出、边际频率与差异度控制：}窗口读出在无理旋转下的 Sturmian 结构、符号 1 的极限频率与有限 $N$ 偏差的差异度界。
  \item \textbf{$X_m$ 计数、$\varphi$-稳定投影与熵率极限：}无相邻 $1$ 语言计数 $|X_m|=F_{m+2}$、投影算子 $P_\varphi$ 的维数与 $\log\varphi$ 熵率。
  \item \textbf{$\Fold_m$ 平均退化度标度与 $\Fold_6$ 显式结构入口：}$d_m=2^m/F_{m+2}$ 的量化表与 $\Fold_6$ 纤维分类的章节入口。
  \item \textbf{$\Fold_6$: 64→21 与退化谱（2/3/4）：}$\Fold_6$ 的 64→21 压缩、退化谱与显式原像列表。
  \item \textbf{$\pi$-通道：循环闭合与 $18\oplus 3$ 分裂：}循环可容许性定义、边界态集合与计数。
  \item \textbf{$\zeta$ 因式分解与极点位置；$e$-通道正规化：}$1-z-z^2$ 的因式分解、极点与归一化说明。
  \item \textbf{Hilbert 编址：$\Sigma_n$ 与 $H_n$ 的定义及局域性要求：}将 tick 索引映射到屏幕格点的编址结构与公理化描述。
  \item \textbf{差异度证书实现与实验表接口：}以黄金斜率实现 $\log 2$、$\pi$ 等平均的实验表与 Koksma 证书接口。
  \item \textbf{Hecke/素数骨架与端到端 toy instance：}cusp 评价、$q$-展开、Hecke 检验与编码链条的端到端实例说明。
  \item \textbf{Zeckendorf 规范化算法与 $\Fold_6$ 计算流程：}贪心 Zeckendorf 算法与 $\Fold_6$ 的实现步骤。
\end{enumerate}

